%!TEX root = ../hutbthesis_main.tex

\chapter{固定翼飞行器着陆控制系统组成与需求分析}

\section{固定翼飞行器着陆任务场景与系统组成}

固定翼飞行器着陆阶段在其整个飞行过程当中，属于风险最为集中且操纵要求颇高的阶段范畴。依据本文仿真平台的运行逻辑剖析可知，飞行器的着陆过程能够被划分成进近、拉平、接地区域、滑跑这四个阶段。进近阶段的主要任务是实现下滑轨迹的维持、速度的有效控制、跑道的精准对准，拉平阶段主要是降低下沉速率并且对俯仰姿态加以调整，让飞行器平稳地向地面靠近，接地区域着重关注接地高度、姿态、侧偏误差，滑跑阶段重点对接地后的方向保持与状态收敛展开分析。着陆后段飞行高度较低、修正时间比较短且对环境扰动比较敏感，所以此阶段既要求控制系统具备稳定可靠性，又要求操纵者能够及时把握关键状态的变化情况并做出正确的决策。

就系统组成而言，本研究的固定翼飞行器着陆控制系统主要涵盖感知层、控制层、人机交互层、安全保护层这四个部分。感知层的作用是收集离地高度、速度、下沉率、俯仰角、滚转角、航向角、侧偏量还有地形高程等状态信息，它是达成闭环控制的根本所在，控制层会按照飞行状态、目标约束，实时产生油门、升降舵、副翼、方向舵、襟翼指令，以此来实现纵向和侧向的修正，人机交互层负责把关键状态、风险提示、建议操纵信息反馈给操纵者，安全保护层会依据阈值判断是不是需要限幅、约束或者复飞。由此能够知道，着陆控制系统并非单纯的自动控制模块，而是一个将状态获取、控制决策、信息提示、保护机制融合在一起的综合系统。

\section{着陆控制关键参数与影响因素}

在固定翼飞行器着陆控制里，关键参数有不少，像目标速度、最小速度与最大速度、下滑角、拉平高度、拉平距离、接地高度、最大下沉率、最大滚转角、最大俯角、最小俯角、最大侧偏量，还有跑道捕获与提交边界等。这里面，速度对飞行器着陆后段有没有足够升力和操纵裕度起着决定作用，下滑角能体现整体下降的趋势，拉平高度和拉平距离决定着阶段切换的时机，下沉率、滚转角、俯仰角、侧偏量会直接对接地品质和安全性产生影响。要是速度低，飞行器操纵也许会变得迟钝，要是下沉率太大，接地冲击会明显增强，要是滚转或者侧偏过大，可能会出现偏离跑道中心线或者接地姿态不正常的情况。

影响着陆控制效果的因素主要可分为三类。其中第一类是环境因素，像地形起伏、局部高程变化、近地扰动等都涵盖在内。当前已有的仿真设置了多种地形，诸如平坦、起伏、山脊、谷地、阶跃等，这表明外部环境发生变化时，飞行器相对地面的离地高度变化规律也会随之直接改变。第二类是飞行器本体因素，主要体现为机体姿态响应能力、起落架缓冲性能。吴磊磊在某无人机油---气式缓冲器强度分析与疲劳寿命预测研究里提到，接地冲击载荷会直接影响缓冲器的受力水平、结构寿命，这也就意味着着陆控制既要关注飞行轨迹，同时也要考虑接地载荷是否处在安全范围之内。第三类是操纵因素，也就是操纵输入是否平稳、修正是否及时、末段判断是否准确。在复杂地形或者近地阶段，如果缺乏有效提示，操纵者容易出现修正滞后或者过度修正的情况，进而影响接地质量。

\section{研究对象参数整理与取值依据}

本文把小型固定翼飞行器着陆控制系统当作研究对象，运用典型参数建模的办法来展开分析，并非针对某一具体实机做完全复现。如此处理，既能凸显控制系统与人机交互强化机制自身，又方便统一不同工况下的仿真比较。按照当前平台的设置，研究对象的初始离地高度是120
ft，初始速度是72 kts，初始下降率为-250 fpm，跑道距离为2200
ft，目标接地点在跑道阈值后120 ft处，跑道捕获距离是320
ft，跑道提交距离是90 ft，跑道提交高度是22
ft。对于近地约束，还设定了flare最大侧偏28
ft、最大航向误差6$^\circ$、最大滚转10$^\circ$、接地提交阶段最大侧偏45
ft、最大航向误差10$^\circ$、最大滚转12$^\circ$。

这些参数的取值依照``典型、合理、可仿真''这一原则。速度参数主要按照小型固定翼飞行器末段进近的通常范围来确定，既要防止失速风险，又得抑制接地冲击，高度和距离参数主要是用来区分进近、拉平和接地区域，以便系统能依照不同阶段调用不同控制策略，横向约束参数是用于保证飞行器在接近目标接地点时具备更严格的跑道对准能力。郑振干在无人机滑橇式起落架动力学特性分析与缓冲性能优化研究里指出，不同着陆构型下的接地响应和冲击分配存在显著差异，所以本文在参数整理时既要考虑飞行控制的需求，又要兼顾接地安全边界与后续分析的统一性。

\section{人机交互强化需求分析}

传统着陆控制方面的研究通常更多地聚焦于飞行器能不能顺着目标轨迹稳定地下降，然而对于操纵者在着陆过程里的感知、判断、干预能力却没有给予足够的考量。事实上，着陆阶段属于一个典型的人机协同进程。结合你现有的平台界面能够看到，系统已经可以实时显示时间、距离、离地的高度、速度、下沉的速率、俯仰角、滚转角、航向角、侧偏量、控制指令，并且能够输出``建议''\,``告警''\,``触发原因''等信息。这说明系统已然具备了进行人机交互强化研究的基础。

本文所持观点为，人机交互强化需求主要呈现于三个方面。其一为状态提示需求，也就是说系统需要及时且清晰地向操纵者展示关键状态量、其变化趋势，以此助力操纵者判定当下所处的飞行阶段、着陆风险。其二是风险告警需求，一旦系统检测出速度过低、下沉率过大、滚转超限、侧偏过大或者跑道对准不足之时，就应当及时发出告警，从而让操纵者能够提前进行修正。其三是操纵辅助与约束需求，即系统在高风险阶段不但要进行提示，而且在必要的时候要通过限幅、约束或者复飞触发来防止误操作使得风险进一步扩大。李月昊于面向无人机起落架的新型惯容式缓冲系统分析与设计研究里提出，借助改善系统动态响应与能量分配路径，能够提高接地适应能力、缓冲效率。这一思路表明，着陆安全的提高不光依赖于结构自身，还依赖于系统对风险的主动识别、过程干预。对于本文来讲，人机交互强化正是经由``提示---辅助---约束''这种方式，提高操纵者在复杂着陆场景下的判断能力与处置能力。本

固定翼飞行器着陆控制系统的需求，不单单是达成基本着陆，而是要在满足轨迹控制、接地安全的条件下，借助人机交互来增强操纵稳定性、风险感知能力、系统整体安全性。这也为后续动力学模型建立、工况设计、交互策略分析打下了基础。
